\section{Additional local-reduction research at the 237-state boundary}

The starting table is the repository's composed 237-state machine, in inherited
labels, with start state 95 and separate halt state 254. Its table hash is
\path{a3c4407f30e2e2139af87c6a91948c518134855841a28eb85b1d81bcfd026c9a}.
The additional searches described here found no sub-237-state candidate.
They extend several previously tested classes, but their negative outputs are
not state lower bounds.

\subsection{Halting observation is weaker than a tape-preserving terminal join}

A relevant omission in the inherited search domains is that replacement
transitions were enumerated over surviving \emph{working} states only. They did
not include halt state 254. Moreover, local joins required equal head positions
and equal tapes even when both computations had halted. Neither restriction is
required by the research objective, which concerns halting from the blank tape.

For the particular 237-state table, the only transition into the halt state is
\[
 \delta(77,1)=(1,+1,254).
\]
Consequently, a fold that deletes state 77, or resynthesizes both its rows while
retaining it, cannot preserve a halting execution if every newly synthesized
transition still excludes target 254. The unique original halting transition
has disappeared and no replacement can enter the halt state. This is a scope
omission in such searches, not an obstruction to merging state 77.

The new terminal variants add 254 to synthesized target domains. They also
identify all pairs of halted configurations, regardless of the final tape or
head position. They retain exact tape and head equality at every nonhalting
checkpoint, after applying the specified state-dependent representation.

\paragraph{Positive-block simulation with terminal observation.}
Let $M$ and $N$ be deterministic Turing machines, let $I$ be an invariant of $M$,
and let $F$ map nonhalting configurations in $I$ to nonhalting configurations of
$N$. Define a relation by
\[
 c\mathrel{R}d\quad\Longleftrightarrow\quad
 \bigl(c\text{ and }d\text{ are both halted}\bigr)
 \ \lor\ 
 \bigl(c\in I\text{ is nonhalting and }d=F(c)\bigr).
\]
Suppose initial configurations are related and that for every nonhalting
$c\in I$ there are positive integers $a(c),b(c)$ such that both finite runs
exist up to their specified endpoints and
\[
 M^{a(c)}(c)\mathrel{R}N^{b(c)}(F(c)).
\]
Then the two initial runs have the same halting status.

\emph{Proof.} Iterating the block condition gives related checkpoints with
strictly increasing clocks in both machines. If $M$ halts after finitely many
steps, the strictly increasing old clock cannot pass the halting time because
a run beyond that time does not exist. It therefore reaches a halted endpoint,
whose related new endpoint is also halted. If $M$ does not halt, the block
condition provides infinitely many new steps because every $b(c)$ is positive;
thus the deterministic new run cannot halt. This proves both implications.
Injectivity of $F$ is not needed.\hfill$\square$

A synthetic control validates this distinction without asserting a new
consistency machine. It has working states $77,95,96$, start 95, and rows
\[
\begin{array}{c|cc}
 &0&1\\\hline
77&(1,L,96)&(1,R,254)\\
95&(1,R,77)&(1,L,96)\\
96&(1,R,77)&(1,R,77)
\end{array}
\]
Merging 96 into 77, redirecting its incoming transition, and setting both rows
of 77 to $(0,L,254)$ yields a two-state machine with the same universal halting
behavior. On the blank tape the old machine takes four steps and halts at head
2 with ones at cells 0 and 1; the new machine takes two steps and halts at head
0 with a one only at cell 0. A terminal relation is therefore necessary for this
particular local comparison. The checker exhausts all 128 assignments on cells
$[-3,3]$ and verifies that neither run reads outside the window, giving a
whole-tape positive control. The search engine independently recovers the
specified merge.

\subsection{Representations at both ends of a proposed merge}

Previous families typically left the retained state's representation unchanged
and modified only the deleted state's representation. Here both states $p,q$
can change their representations. Suppose every external incoming transition
to $s\in\{p,q\}$ has the same direction $d_s\in\{-1,+1\}$. An external
predecessor is a state other than $p$ and $q$. Choose $\epsilon_s\in\{0,1\}$
and one of the four Boolean functions
\[
 f_s(x)\in\{x,0,1,1-x\}.
\]
The proposed representation at $s$ is
\[
 F(s,h,t)=
 \left(p,h-2\epsilon_s d_s,
 t[h-d_s\leftarrow f_s(t(h-d_s))]\right).
\]
All other nonhalting states initially have the identity representation.
The new tape may keep, erase, set, or negate the cell written by the caller;
its head may stay aligned with the old head or shift by two cells.

For an external incoming transition
\[
 \delta(r,b)=(w,d_s,s),\qquad r\notin\{p,q\},
\]
the required replacement is exactly
\[
 \delta'(r,b)=\bigl(f_s(w),(-1)^{\epsilon_s}d_s,p\bigr).
\]
This formula has a direct proof on every tape. If the old head at the caller is
$k$, the old successor's head is $k+d_s$. Applying $F$ places the new head at
$k+d_s-2\epsilon_s d_s=k+(-1)^{\epsilon_s}d_s$ and changes only cell $k$, to
$f_s(w)$. This is precisely the replacement transition. The statement remains
valid when $f_s$ is constant; preservation of halting does not require an
injective representation.

The two rows at the merged state are synthesized freely. Their compatibility
must be checked, because the caller formula alone proves no correctness claim
for the merge. A further family also changes one helper row from the inherited
neutral catalogue and rechecks its original contexts together with the merged
states' contexts. Individual neutrality is never assumed to imply compatibility
with a changed representation.

The strict common-row search covers 180 eligible state pairs and 49 nonidentity
encoding pairs per state pair, for 8,820 representations. The terminal-aware
search focused on the halting gate covers 207 eligible pairs involving 77 and
all 64 encoding pairs, including the identity, for 13,248 representations.
Both require a common external incoming direction at each encoded label and
exclude encoding the initial state 95. These restrictions leave many possible
representations and transformations outside the search.

\subsection{Freely synthesized helper compensation}

Another new search drops the requirement that each helper edit be independently
neutral in the original machine. It considers every pair sharing one literal
row, preserves that common row, and synthesizes the other retained row. It
chooses a first helper edit from the full finite transition domain. The first
helper must make all merged-state profile cases pass. If its original contexts
then fail, a second helper edit is synthesized along the first failing new
trace. The final check includes the merged states and the original contexts of
both helpers.

This is a sequential compensation subclass. It excludes possibilities where
both helper edits are needed before the first helper can pass the merged-state
cases. It also excludes arbitrary larger clusters and unrestricted changes to
callers. At horizon 17, all 234 common-row pairs were completed. The two runs
partitioning those pairs tried a total of 220,896 retained-row choices,
1,491,793,070 first-helper choices and 16,878,220 second-helper choices, with no
candidate. A terminal-aware follow-up at horizon 65 covered the two common-row
pairs involving 77, namely $77/204$ and $77/252$, and likewise found none.

\subsection{Evidence, controls, and limits}

Discovery uses the retained source-constrained finite profile of 29,192
state/window pairs, with window $[-5,5]$. Old and new traces must remain within
the known window while reading or writing. Both clocks have the specified
finite horizon, and the new trace may join any positive old checkpoint within
that horizon. These finite cases are not an exhaustive description of the
source runtime. An independent check confirms that every profile window belongs
to the inductively closed all-tape sliding-window abstraction; the latter has
93,177 state/window pairs including the halt state. This inclusion does not
make the smaller profile exhaustive.

The original helper catalogue contains 601 finite neutral substitutions.
Regeneration with direct halt targets and the terminal relation considers
450,774 substitutions and returns 604. The only additions are
\[
 77/1\mapsto0L254,\qquad77/1\mapsto0R254,\qquad77/1\mapsto1L254.
\]
Every helper search focused on a merge involving 77 excludes 77 itself from
helper selection. The active helper domains are therefore exactly identical
for the 601- and 604-entry catalogues. The package checks this equality directly;
there is no need to repeat an identical search after changing the catalogue.

A separate positive control recovers the previously established $235\to98$
fold from the 238-state table. The independent checker reconstructs the exact
237-state table, verifies the predecessor invariant in 30,112 Boolean-window
cases, checks the three open-tape exceptional blocks of lengths $3,5,3$, and
checks 473 ordinary commuting old rows. These are sufficient finite premises
for an ordinary all-tape proof of that inherited fold. The control does not
claim a new reduction.

The supplied JSON receipt and logs give commands, input and code hashes,
completed counters, horizon choices, and positive-control outputs. No compiled
binaries are needed. The retained discoveries are negative within these stated
classes; there is no proof-assistant theorem, no new 236-state table, no
minimality claim, and no assertion that broader coordinated edits cannot work.
